% Lösungen Rekonstruktion von Funktionsgleichungen (zusammenbau v0.4, Heft)
\einheitenkopf{Rekonstruktion von Funktionsgleichungen}

\erg{14}{a) genau passend – drei Unbekannte $a$, $b$, $c$ und drei Gleichungen. \quad b) $m = 3$, $n = 2$, also $y = 3x + 2$ \quad c) Ansatz $f(x) = ax^4 + bx^2 + c$: $c = 4$; $a + b = -3$ und $16a + 4b = 0$; $a = 1$, $b = -4$, also $f(x) = x^4 - 4x^2 + 4$}
\erg{15}{a) $h(0{,}6) = 1{,}2$ und $h(3) = 6$ – beide Punkte liegen auf dem Graphen, und durch zwei Punkte geht genau eine Gerade. \quad b) $k(-1) = 0{,}5 + 2 = 2{,}5$ und $k(3) = -1{,}5 + 2 = 0{,}5$ – beide Punkte liegen auf der Geraden, und durch zwei Punkte geht genau eine Gerade. \quad c) $s(1{,}5) = 3 - 4 = -1$ und $s(3{,}5) = 7 - 4 = 3$ – beide Punkte liegen auf dem Graphen, und durch zwei Punkte geht genau eine Gerade. \quad d) $f(0) = a = 2$; $f(1) = 2b = 6$, also $b = 3$}
\erg{16}{a) zwei Gleichungen – $f(3) = 1$ und $f'(3) = 0$. \quad b) $c = -1$; $f(1) = a + b - 1 = -3$ und $f'(1) = 2a + b = 0$; $a = 2$, $b = -4$, also $f(x) = 2x^2 - 4x - 1$ \quad c) $c = 4$; $f(-1) = -4 + 7 = 3$ und $f'(-1) = 4$: $a - b + 4 = 3$ und $-2a + b = 4$; $a = -3$, $b = -2$, also $f(x) = -3x^2 - 2x + 4$}
\erg{17}{a) Ansatz $q(x) = ax^2 + bx$ wegen $q(0) = 0$; $q(3) = 9a + 3b = 2{,}7$ und $q'(3) = 6a + b = 0$; $a = -0{,}3$, $b = 1{,}8$, also $q(x) = -0{,}3x^2 + 1{,}8x$ \quad b) $h'(t) = 3at^2 + 2bt$; $h(4) = 64a + 16b = 64$ und $h'(4) = 48a + 8b = 24$; $a = -0{,}5$, $b = 6$, also $h(t) = -0{,}5t^3 + 6t^2$ \quad c) Aus I und II: $d = 15$, $c = 4$; $w(4) = 64a + 16b + 31 = 31$ und $w'(4) = 48a + 8b + 4 = 0$; $a = -0{,}25$, $b = 1$, also $w(x) = -0{,}25x^3 + x^2 + 4x + 15$ \quad d) $c = 0$; $f(1) = 3 \cdot 1 + 1 = 4$ und $f'(1) = 3$: $a + b = 4$ und $2a + b = 3$; $a = -1$, $b = 5$, also $f(x) = -x^2 + 5x$ \quad e) $c = 0$; $f(3) = -2 \cdot 3 + 9 = 3$ und $f'(3) = -2$: $9a + 3b = 3$ und $6a + b = -2$; $a = -1$, $b = 4$, also $f(x) = -x^2 + 4x$ \quad f) Knickfrei in $S(0|2{,}5)$: $f(0) = r(0) = 2{,}5$ und $f'(0) = r'(0) = 0$, also $c = 2{,}5$, $b = 0$; $f(4) - g(4) = 1{,}2$ mit $g(4) = 0{,}5$: $16a + 2{,}5 = 1{,}7$, also $a = -0{,}05$ und $f(x) = -0{,}05x^2 + 2{,}5$ \quad g) $k(1) = a + b = 3$ und $k'(1) = 4a + 2b = 8$; $a = 1$, $b = 2$}
\erg{18}{a) $w(5) = 20 + b \cdot e^{5c} = 80$, also $b \cdot e^{5c} = 60$; $w'(5) = b \cdot c \cdot e^{5c} = 60c = -6$, also $c = -0{,}1$ und $b = 60 \cdot e^{0{,}5} \approx 98{,}92$ \quad b) $w(2) = 22 + b \cdot e^{2c} = 7$, also $b \cdot e^{2c} = -15$; $w'(2) = b \cdot c \cdot e^{2c} = -15c = 3$, also $c = -0{,}2$ und $b = -15 \cdot e^{0{,}4} \approx -22{,}38$ \quad c) $w(4) = 12 + b \cdot e^{4c} = 18$, also $b \cdot e^{4c} = 6$; $w'(4) = b \cdot c \cdot e^{4c} = 6c = -1{,}5$, also $c = -0{,}25$ und $b = 6 \cdot e^{1} \approx 16{,}31$}
\erg{19}{a) $6$ – vom Hoch- zum Tiefpunkt ist es eine halbe Periode, zum nächsten Hochpunkt eine ganze. \quad b) $a = \frac{3 - (-3)}{2} = 3$; halbe Periode $3 - 1 = 2$, also $p = 4$ und $b = \frac{2\pi}{4} = \frac{\pi}{2} \approx 1{,}57$ \quad c) Viertelperiode $3{,}1$: $p = 12{,}4$ und $b = \frac{2\pi}{12{,}4} \approx 0{,}51$; $a = 1{,}8$}
\erg{20}{a) Viertelperiode $3$: $p = 12$ und $b = \frac{2\pi}{12} = \frac{\pi}{6} \approx 0{,}52$; $a = g(3) = f(3) = 2$ \quad b) Viertelperiode $4$: $p = 16$ und $b = \frac{2\pi}{16} = \frac{\pi}{8} \approx 0{,}39$; $a = g(4) = f(4) = 4$}
\erg{21}{a) $p(0) = 3$, $p'(0) = -2$, $p(1) = 0{,}5$, $p'(1) = 1$; aus den ersten drei: $c = 3$, $b = -2$, $a = -0{,}5$; Probe: $p'(1) = 2a + b = -3$ – das ist nicht $1$, es gibt keine solche Funktion. \quad b) $p(0) = -1$, $p'(0) = 1$, $p(1) = 1$, $p'(1) = 3$; aus den ersten drei: $c = -1$, $b = 1$, $a = 1$; Probe: $p'(1) = 2a + b = 3$ – die vierte Bedingung ist erfüllt, es gibt $p$ mit $p(x) = x^{2} + x - 1$. \quad c) $p(0) = 4$, $p'(0) = 0$, $p(2) = 0$, $p'(2) = -3$; aus den ersten drei: $c = 4$, $b = 0$, $a = -1$; Probe: $p'(2) = 4a + b = -4$ – das ist nicht $-3$, es gibt keine solche Funktion.}
